% CS F407 Artificial Intelligence: final project report template.
% Limit: 2 pages, including everything except references.
% Do not change the font size, margins or spacing.
\documentclass[11pt,a4paper,twocolumn]{article}

\usepackage[margin=2cm,columnsep=0.8cm]{geometry}
\usepackage[T1]{fontenc}
\usepackage{lmodern}
\usepackage{microtype}
\usepackage{amsmath,amssymb}
\usepackage{graphicx}
\usepackage{booktabs}
\usepackage{tabularx}
\usepackage{enumitem}
\usepackage{titlesec}
\usepackage[hidelinks]{hyperref}

\titleformat{\section}{\normalsize\bfseries}{\thesection}{0.6em}{}
\titlespacing*{\section}{0pt}{1.2ex plus 0.3ex}{0.6ex}
\setlist{nosep,leftmargin=*}
\setlength{\parskip}{0.3ex}

\newcommand{\coursename}{CS F407 Artificial Intelligence, AY 2026--27 Semester I}
\newcommand{\teamno}{Team [TODO]}
\newcommand{\githuburl}{https://github.com/krishavgoyal/Cross_Era_Fantasy_Cricket}

\begin{document}

\twocolumn[{%
\begin{center}
{\large\bfseries Cross-Era Fantasy Cricket: Animating Historical Cricketers with Modern Motion}\\[1ex]
\teamno\\
Pranahv Subramaniam Sankar (2024A4PS0423G), Ahaan Pathak (2024A3PS0396G),\\ Ankit Pandey (2024AAPS0808G), Krishav Goyal (2024A7PS0855G)\\[0.5ex]
{\small \coursename, BITS Pilani, K.K. Birla Goa Campus}\\[1.5ex]
\begin{minipage}{0.85\textwidth}
\small
\centerline{\bfseries Abstract}\vspace{0.5ex}
Fantasy cricket compares players across eras only through statistics. We make such matchups visible: from one
archival photograph and a modern clip we generate a short video of the historical player performing the modern
motion, styled like footage of his era. DWPose extracts the motion on a laptop CPU; Wan2.1-VACE-1.3B generates the
clip on a free Kaggle T4 from the photo and a control video; post-processing retimes it and applies 1930s/1970s/1990s
filters. In a controlled ablation, a depth control instead of a pose skeleton kept the bat in 5/5 runs (0/2 with the
skeleton) and lowered brightness flicker from 2.38 to 1.68, while motion accuracy stayed high (PCK@0.2 0.90--1.00).
A skeleton-only baseline (MusePose) shows why fast bowling is harder: blurred source frames and front/back ambiguity.
\end{minipage}
\end{center}
\vspace{1.5ex}
}]

\section{Introduction}
Sir Don Bradman never faced a modern fast bowler; our system shows what it could look like. \textbf{Input:} one
photograph of a historical player (appearance) and a clip of a modern player (motion). \textbf{Output:} a 2--6\,s video
of the historical player performing that motion in the visual style of his era, labelled ``AI-generated''. We
contribute (i) a free, reproducible pipeline (laptop CPU and a free Kaggle GPU); (ii) a player-only depth control that
keeps the bat and removes bystanders; (iii) a controlled ablation with several seeds; and (iv) an analysis of why a
skeleton-driven baseline fails on fast bowling.

\section{Background and related work}
The task is reference-guided human image animation with conditional generative models. Pose estimators such as
DWPose~\cite{yang2023dwpose} give 2D joint sequences. Diffusion models with a reference branch, such as Animate
Anyone~\cite{hu2023animateanyone}, MagicAnimate~\cite{xu2023magicanimate} and MusePose~\cite{musepose2024}, animate a
photo from a pose sequence; AnimateDiff~\cite{guo2023animatediff} adds temporal layers to image diffusion. Video
diffusion transformers such as Wan2.1~\cite{wan2025} with VACE~\cite{jiang2025vace} accept a reference image and a pose
or depth video. Depth Anything V2~\cite{yang2024depthanythingv2} estimates monocular depth. The proposal named
LTX-2~\cite{ltx2}, which uses the photo only as a first frame and needs 32\,GB+ GPUs, so we chose VACE.

\section{Method}
\textbf{Motion.} rtmlib DWPose-l runs on CPU; the player is chosen by a box on frame 0 and followed by best-IoU
tracking. Joints with confidence $<0.3$ are masked, gaps of up to 4 frames interpolated, tracks smoothed
(Savitzky--Golay, window 7) and resampled to 33 frames on a $480{\times}832$ canvas with an 8\% top margin.
\textbf{Controls.} \emph{Skeleton}: an OpenPose-style render. \emph{Depth}: Depth Anything V2 on the same frames and
crop, with one global normalisation (no flicker) and the head flattened so a modern helmet does not replace the face.
\emph{Player-only depth} (bowler): a person mask kept only where it contains the player's keypoints, so fielders vanish.
\textbf{Reference.} rembg removes the photo background; the player is tightly cropped on white.
\textbf{Generation.} Wan2.1-VACE-1.3B (fp16) in a headless ComfyUI on a Kaggle T4: 20 steps, CFG 5, shift 8. The prompt
describes 1930s kit; the negative prompt lists observed failures (extra balls, two bats, logos, hats).
\textbf{Post-processing.} Frame blending retimes to real speed at 24\,fps; seeded era filters (grain, flicker and gate
weave for the 1930s; scanlines and colour fringing for the 1970s; chroma bleed for the 1990s) burn in the label.

\section{Experimental setup}
\textbf{Data.} A modern cover drive (5.6\,s) and a front-on fast bowler (2.0\,s) supply motion; public-domain photos of
Bradman (batting) and Harold Larwood (bowling, 1929) supply appearance.
\textbf{Runs.} A $2{\times}2$ ablation, control (skeleton/depth) $\times$ reference (padded/tight crop), 17 frames;
seeds 42/7/123/2024 for depth with the tight crop; final 33-frame runs for both players.
\textbf{Metrics.} PCK@0.2: DWPose on the output vs the control joints (17 joints, threshold $0.2\times$ torso length).
Identity: DINOv2 cosine similarity of person crops vs the photo. Flicker: std of per-frame brightness. Bat: one rigid
bat in every frame, by inspection.
\textbf{Hardware.} A laptop RTX 4050 (6\,GB) for everything except generation; on the Kaggle T4 a 33-frame clip takes
${\approx}520$\,s at a 14.9\,GB peak, so 33 frames is the limit.
\textbf{Baseline.} Skeleton control, and the team's skeleton-driven MusePose pipeline (Section~6).

\section{Results}
\begin{table}[h]
\centering
\small
\setlength{\tabcolsep}{3pt}
\begin{tabular}{lcccc}
\toprule
Control (runs) & PCK & Identity & Flicker & Bat\\
\midrule
Skeleton (A, C) & 0.95 & 0.19 & 2.38 & 0/2\\
Depth (B, D1--D4) & 0.97 & 0.23 & 1.68 & 5/5\\
\midrule
Depth, final batter & 0.90 & 0.22 & 1.65 & yes\\
Player-only depth, bowler & 0.81 & 0.11 & 3.67 & --\\
\bottomrule
\end{tabular}
\caption{Main results (means over runs). PCK@0.2 and identity: higher is better; flicker: lower is better; Bat: runs
that keep one bat throughout.}
\label{tab:results}
\end{table}
\textbf{Depth fixes the bat.} With only a skeleton, which never describes the bat, the model loses or bends it; with
depth, one rigid bat is held in every frame of all five runs, with no loss of motion accuracy and less flicker
(Table~\ref{tab:results}). \textbf{The tight crop did not help likeness}: with the seed fixed it changed identity by
$-0.02$ to $-0.08$, while the seed alone moved it by up to 0.14, so likeness comparisons need several seeds.
\textbf{Bowler.} Player-only depth removed the fielders, and the negative prompt removed the 1929 photo's fedora and
suit. PCK is lower (0.79--0.85) because the bowling arm is the fastest, most blurred motion.
\textbf{Era filters} keep the motion: PCK falls by at most 0.08 (batter) and 0.12 (bowler).

\section{Discussion and limitations}
\textbf{Why the skeleton baseline failed on bowling.} The baseline fits a rigid 3D skeleton to DWPose and MediaPipe
detections (fixed bone lengths, planted feet, rigid bat) and drives MusePose. It worked for batting (Bradman PCK@0.1
0.963, merged limbs 0.7\%) but not for bowling (McGrath PCK@0.1 0.64). Its bowling source was a dark, side-on
slow-motion recording in which the arm is blurred and the feet leave the frame, so those joints were guessed. A
skeleton also does not say which limb is in front: the generated bowler was detected mirrored in up to 60\% of frames
(allowing mirroring, PCK is 0.88), and the reference photo faced the opposite way to the motion. The VACE bowler used a
sharp, front-on, fully visible clip, a depth control that fixes front/back order, and a front-facing photo.
\textbf{Limitations.} Face likeness is weak (identity 0.1--0.3), the helmet outline can become a dark cap, clips are
short (the T4 limit) and the two players are generated separately. \textbf{Next:} combine depth and skeleton control,
longer clips on a larger GPU, and a face-identity adapter.

\section{Conclusion}
A free pipeline animates historical cricketers convincingly for batting. The control signal matters more than the
reference crop: depth keeps the bat and the body's front/back order, and clean source footage matters most for fast
bowling.

\section*{Individual member contributions}
{\small
\begin{tabularx}{\linewidth}{@{}l X@{}}
\toprule
Name & Contribution \\
\midrule
Ankit Pandey & Problem formulation, pipeline overview, limitations and ethics, Sections 1, 2 and 7. \\
Pranahv S. Sankar & Motion extraction (DWPose tracking, skeleton and depth controls), rigid 3D skeleton fit, integration and the MusePose baseline, Sections 3 and 6. \\
Ahaan Pathak & Model selection and running the model (VACE on a headless Kaggle T4), AnimateDiff/MagicAnimate team baseline, Section 4. \\
Krishav Goyal & Generation and ablation, era styling, evaluation metrics and results, Section 5. \\
\bottomrule
\end{tabularx}}

\smallskip
\noindent All members have read this report and agree with the statement above.

\section*{Generative AI disclosure}
We used Claude (Anthropic): the Claude chat app for background research and step-by-step guidance, and Claude Code
(Claude Opus 5.5) to write and debug the pipeline scripts, the Kaggle notebook and the 3D skeleton fit, and to draft
and edit the report text. It affected the code in \texttt{scripts/}, \texttt{kaggle/}, \texttt{motion3d/} and
\texttt{animate/}, and all sections of this report. We checked all such output, ran the code ourselves, verified every
reported number against our own runs, and take responsibility for it. All generated videos carry an
``AI-generated -- academic project'' label.

\section*{Code availability}
\noindent Code and data:\\
\url{\githuburl}\\
The repository is public and its README gives instructions to reproduce the results reported here.

\begin{thebibliography}{9}
\small
\bibitem{yang2023dwpose}
Z.~Yang, A.~Zeng, C.~Yuan and Y.~Li.
\newblock Effective whole-body pose estimation with two-stages distillation.
\newblock \emph{arXiv:2307.15880}, 2023.
\bibitem{hu2023animateanyone}
L.~Hu, X.~Gao, P.~Zhang, K.~Sun, B.~Zhang and L.~Bo.
\newblock Animate Anyone: consistent and controllable image-to-video synthesis for character animation.
\newblock \emph{arXiv:2311.17117}, 2023.
\bibitem{xu2023magicanimate}
Z.~Xu, J.~Zhang, J.~H.~Liew, H.~Yan, J.-W.~Liu, C.~Zhang, J.~Feng and M.~Z.~Shou.
\newblock MagicAnimate: temporally consistent human image animation using diffusion model.
\newblock \emph{arXiv:2311.16498}, 2023.
\bibitem{musepose2024}
Z.~Tong, C.~Li, Z.~Chen, B.~Wu and W.~Zhou.
\newblock MusePose: a pose-driven image-to-video framework for virtual human generation.
\newblock \url{https://github.com/TMElyralab/MusePose}, 2024.
\bibitem{guo2023animatediff}
Y.~Guo, C.~Yang, A.~Rao, Z.~Liang, Y.~Wang, Y.~Qiao, M.~Agrawala, D.~Lin and B.~Dai.
\newblock AnimateDiff: animate your personalized text-to-image diffusion models without specific tuning.
\newblock \emph{arXiv:2307.04725}, 2023.
\bibitem{wan2025}
Wan Team.
\newblock Wan: open and advanced large-scale video generative models.
\newblock \emph{arXiv:2503.20314}, 2025.
\bibitem{jiang2025vace}
Z.~Jiang, Z.~Han, C.~Mao, J.~Zhang, Y.~Pan and Y.~Liu.
\newblock VACE: all-in-one video creation and editing.
\newblock \emph{arXiv:2503.07598}, 2025.
\bibitem{yang2024depthanythingv2}
L.~Yang, B.~Kang, Z.~Huang, Z.~Zhao, X.~Xu, J.~Feng and H.~Zhao.
\newblock Depth Anything V2.
\newblock \emph{arXiv:2406.09414}, 2024.
\bibitem{ltx2}
Lightricks.
\newblock LTX-2: efficient joint audio-visual foundation model.
\newblock \emph{arXiv:2601.03233}, 2026.
\end{thebibliography}

\end{document}
